% Lean: EconLean/GE/Production.lean
\chapter{Production Economies and the Welfare Theorems}
\label{chap:ge_production}

Source: MWG \S16.B--16.D. A \emph{private ownership economy} adds firms $j$ with production sets
$Y_j \subseteq \R^\iota$ to an exchange economy (Definition~\ref{def:exchange_economy}); consumer $a$ owns share
$\theta_{aj} \ge 0$ of firm $j$, with $\sum_a \theta_{aj} = 1$.

\begin{definition}[Private ownership economy; MWG \S16.B]
\label{def:production_economy}
\lean{Econ.ProductionEconomy, Econ.ProductionEconomy.ProfitMax, Econ.ProductionEconomy.wealth,
  Econ.ProductionEconomy.Feasible, Econ.ProductionEconomy.ParetoOptimal,
  Econ.ProductionEconomy.IsWalrasianEquilibrium,
  Econ.ProductionEconomy.IsPriceQuasiEquilibriumWithTransfers,
  Econ.ProductionEconomy.IsPriceEquilibriumWithTransfers}
\leanok
\uses{def:exchange_economy}
$y$ \emph{maximizes profit} on $Y$ at $p$ if $y \in Y$ and $p \cdot y' \le p \cdot y$ for $y' \in Y$. Consumer $a$'s
wealth is $p \cdot \omega_a + \sum_j \theta_{aj}\, p \cdot y_j$. $(x, y)$ is \emph{feasible} if $x_a \ge 0$,
$y_j \in Y_j$ and $\sum_a x_a = \bar\omega + \sum_j y_j$; \emph{Pareto optimal} if feasible and no feasible $(x', y')$
Pareto dominates $x$. $(p, x, y)$ is a \emph{Walrasian equilibrium} if every firm maximizes profit, every $x_a$ is
demanded at that wealth, and markets clear. A \emph{price (quasi-)equilibrium with transfers} $(p, w, x, y)$ has
$(x,y)$ feasible, firms maximizing profit, $\sum_a w_a = p \cdot (\bar\omega + \sum_j y_j)$, and every $z \ge 0$
with $z \spref_a x_a$ costing more than (at least) $w_a$.
\end{definition}

\begin{lemma}[Total wealth]
\label{lem:sum_wealth}
\lean{Econ.ProductionEconomy.sum_wealth}
\leanok
\uses{def:production_economy, lem:cost_sum}
$\sum_a \mathrm{wealth}_a(p, y) = p \cdot (\bar\omega + \sum_j y_j)$.
\end{lemma}
\begin{proof}
\leanok
Split the sum, use linearity of $p \cdot$ (Lemma~\ref{lem:cost_sum}), swap the order of summation
(\texttt{Finset.sum\_comm}), and use $\sum_a \theta_{aj} = 1$.
\end{proof}

\begin{theorem}[First Welfare Theorem with production; MWG Prop.~16.C.1]
\label{thm:first_welfare_production}
\lean{Econ.ProductionEconomy.first_welfare_theorem}
\leanok
\uses{def:production_economy}
If every $\pref_a$ is rational and locally nonsatiated, every Walrasian equilibrium allocation $(x, y)$ is Pareto
optimal. No assumption on the production sets is needed.
\end{theorem}
\begin{proof}
\leanok
\uses{lem:cost_ge_demand, lem:cost_gt_demand, lem:sum_wealth}
An equilibrium allocation is feasible. Suppose a feasible $(x', y')$ Pareto dominates $x$, with $b$ strictly better
off. By Lemma~\ref{lem:cost_ge_demand} $p \cdot x'_a \ge \mathrm{wealth}_a$ for all $a$, and by
Lemma~\ref{lem:cost_gt_demand} the inequality is strict for $b$ (\texttt{Finset.sum\_lt\_sum}). With
Lemma~\ref{lem:sum_wealth}: $p \cdot (\bar\omega + \sum_j y_j) < \sum_a p \cdot x'_a = p\cdot(\bar\omega + \sum_j y'_j)$.
Profit maximization gives $p \cdot y'_j \le p \cdot y_j$ for each $j$; summing, a contradiction.
\end{proof}

\begin{lemma}[Limit lemma]
\label{lem:le_sum_cost_closure}
\lean{Econ.ProductionEconomy.le_sum_cost_of_mem_closure, Econ.ProductionEconomy.sum_update_eq}
\leanok
If $c_0 \le \sum_a p \cdot y'_a$ whenever $y'_a \in s_a$ for all $a$, then $c_0 \le \sum_a p \cdot z_a$ whenever
$z_a \in \overline{s_a}$ for all $a$. Also $\sum_b f(x[a \mapsto v]_b) = \sum_b f(x_b) - f(x_a) + f(v)$.
\end{lemma}
\begin{proof}
\leanok
The map $y' \mapsto \sum_a p \cdot y'_a$ is continuous, so $\{y' : c_0 \le \sum_a p\cdot y'_a\}$ is closed and contains
the product $\prod_a s_a$; its closure is $\prod_a \overline{s_a}$ (\texttt{closure\_pi\_set}). The update identity is
\texttt{Finset.sum\_update\_of\_mem} with \texttt{Finset.add\_sum\_erase}.
\end{proof}

\begin{theorem}[Second Welfare Theorem with production; MWG Prop.~16.D.1]
\label{thm:second_welfare_production}
\lean{Econ.ProductionEconomy.second_welfare_theorem}
\leanok
\uses{def:production_economy}
Let $\iota$ and the set of consumers be nonempty, every $\pref_a$ rational, convex and locally nonsatiated, and every
$Y_j$ convex. For every Pareto optimal $(x, y)$ there is $p \neq 0$ such that $(p, w, x, y)$ with $w_a = p \cdot x_a$ is
a price quasi-equilibrium with transfers (in particular every firm maximizes profit at $p$).
\end{theorem}
\begin{proof}
\leanok
\uses{lem:strictUpper, thm:price_support, lem:le_sum_cost_closure}
Let $V_a = \{z \ge 0 : z \spref_a x_a\}$, $V = \{\sum_a x'_a : x'_a \in V_a\}$ and
$Z = \{\bar\omega + \sum_j y'_j : y'_j \in Y_j\}$. Both are convex (Lemma~\ref{lem:strictUpper}(ii); convexity of $Y_j$),
hence so is $V - Z$ (\texttt{Convex.sub}). $0 \notin V - Z$, since otherwise $(x', y')$ would be a feasible allocation
making every consumer strictly better off, contradicting Pareto optimality (there is a consumer $a_0$). So
$0 \notin \operatorname{int}(V - Z)$ and Theorem~\ref{thm:price_support} gives $p \neq 0$ with
$p \cdot z \le p \cdot v$ for $v \in V$, $z \in Z$. By feasibility $\sum_a p \cdot x_a = p \cdot(\bar\omega + \sum_j y_j)$,
which is the transfer condition.

\emph{Firms.} Fix $j$ and $y' \in Y_j$. Then $\bar\omega + \sum_k y[j\mapsto y']_k \in Z$, so its value is
$\le \sum_a p \cdot x'_a$ for all $x'$ with $x'_a \in V_a$; by Lemma~\ref{lem:le_sum_cost_closure} and
$x_a \in \overline{V_a}$ (Lemma~\ref{lem:strictUpper}(iii)) it is $\le \sum_a p \cdot x_a = p\cdot(\bar\omega + \sum_k y_k)$.
Expanding with the update identity gives $p \cdot y' \le p \cdot y_j$.

\emph{Consumers.} Fix $a$ and $z \ge 0$ with $z \spref_a x_a$. Apply Lemma~\ref{lem:le_sum_cost_closure} with
$s_a = \{z\}$ and $s_b = V_b$ for $b \ne a$, benchmark $\bar\omega + \sum_j y_j \in Z$, at the point $x[a \mapsto z]$:
$\sum_b p\cdot x_b \le \sum_b p \cdot x_b - p\cdot x_a + p \cdot z$, i.e.\ $p\cdot x_a \le p \cdot z$.
\end{proof}

\begin{theorem}[Quasi-equilibrium to equilibrium with production; MWG Prop.~16.D.2]
\label{thm:quasi_to_eq_production}
\lean{Econ.ProductionEconomy.IsPriceQuasiEquilibriumWithTransfers.isPriceEquilibriumWithTransfers}
\leanok
\uses{def:production_economy, def:continuousPref}
With rational, continuous preferences and a cheaper point $z \ge 0$, $p \cdot z < w_a$, for every consumer, a price
quasi-equilibrium with transfers is a price equilibrium with transfers.
\end{theorem}
\begin{proof}
\leanok
\uses{lem:cheaper_point}
Feasibility, profit maximization and the transfer condition carry over; the consumer condition is
Lemma~\ref{lem:cheaper_point} applied to each consumer.
\end{proof}

\section{Existence of equilibrium with production}
\label{sec:existence_production}

\begin{lemma}[Demand blows up at boundary prices]
\label{lem:demand_unbounded}
\lean{Econ.RegularUtility.demand_unbounded, Econ.exists_coord_bound}
\leanok
\uses{def:regular_utility}
Let $u$ be regular. If $q_n \gg 0$ converge to $\bar p$ with some $\bar p_l = 0$, wealth $w_n$ stays above some $c > 0$, and
$x_n$ is demanded at $(q_n, w_n)$, then $\max_k x_{n,k} \to \infty$.
\end{lemma}
\begin{proof}
\leanok
Otherwise a subsequence of $x_n$ stays in a box and converges to $\bar x \ge 0$. Since $w_n \ge \max(c, q_n \cdot x_n) \to
W = \max(c, \bar p \cdot \bar x) > 0$ and $\bar p \cdot \bar x \le W$, the improvement property gives $y$ with $u(y) > u(\bar x)$
and $\bar p \cdot y < W$. Then $y$ is eventually affordable, so $u(y) \le u(x_n) \to u(\bar x)$, a contradiction.
\end{proof}

\begin{theorem}[Existence with production]
\label{thm:existence_production}
\lean{Econ.ProductionEconomy.exists_walrasianEquilibrium_production,
  Econ.ProductionEconomy.exists_continuous_supply,
  Econ.ProductionEconomy.exists_walrasianEquilibrium_production_of_unique}
\leanok
\uses{lem:demand_unbounded, thm:existence_regular, lem:sum_wealth}
Suppose consumers have regular utilities, $\bar\omega \gg 0$ (individual endowments may be zero), and firms have compact production
sets containing $0$ with a profit-maximizing supply function continuous on $p \gg 0$ (for instance, unique profit
maximizers, by Berge's theorem). Then a Walrasian equilibrium with $p \gg 0$ exists. Compactness stands in for
Arrow--Debreu's closed, convex, free-disposal sets (it excludes free disposal and constant returns).
\end{theorem}
\begin{proof}
\leanok
Set $w_a(p) = p\cdot\omega_a + \sum_j \theta_{aj}\, p \cdot s_j(p) \ge p \cdot \omega_a \ge 0$ (profits are nonnegative as
$0 \in Y_j$); at zero wealth demand is $0$, and it is continuous there since $x_i \le w/p_i$ and $z(p) = \sum_a x_a(p, w_a(p)) - \bar\omega - \sum_j s_j(p)$. It is continuous on $p \gg 0$, satisfies
Walras' law, and is bounded below since production sets are bounded. Near a boundary price some consumer has
$\bar p \cdot \omega_a > 0$, so their wealth stays bounded away from $0$ and Lemma~\ref{lem:demand_unbounded} gives the
boundary condition. The interior existence theorem gives $z(p) = 0$.
\end{proof}

\begin{theorem}[A disc technology]
\label{thm:disc_economy}
\lean{Econ.GEExamples.norm2, Econ.GEExamples.discSet, Econ.GEExamples.discSupply, Econ.GEExamples.norm2_pos,
  Econ.GEExamples.norm2_sq, Econ.GEExamples.zero_mem_discSet, Econ.GEExamples.discSet_noFreeLunch, Econ.GEExamples.isCompact_discSet,
  Econ.GEExamples.discSupply_profitMax, Econ.GEExamples.continuousOn_discSupply, Econ.GEExamples.discEconomy,
  Econ.GEExamples.discEconomy_exists_equilibrium, Econ.GEExamples.discEconomy3,
  Econ.GEExamples.discEconomy3_endow_two, Econ.GEExamples.discEconomy3_exists_equilibrium}
\leanok
\uses{thm:existence_production, thm:sqrt_economy}
A firm with production set the Euclidean disc of centre $(-1, -2)$ and radius $\sqrt5$ (so $0$ is on its boundary, and
the disc lies in $y_0 + 2y_1 \le 0$, so no nonzero nonnegative output is possible) has supply $c + \rho\, p/\|p\|_2$. Two square-root consumers owning $(1,0)$ and $(0,1)$ and half the firm each
form an economy with a Walrasian equilibrium; so does a variant with a third consumer who owns no goods and the whole
firm.
\end{theorem}
\begin{proof}
\leanok
Profit maximization is Cauchy--Schwarz: $p\cdot(y - c) \le \|p\|_2 \|y - c\|_2 \le \|p\|_2\,\rho$. Apply
Theorem~\ref{thm:existence_production}.
\end{proof}

\section{Existence with unbounded production sets}
\label{sec:existence_unbounded}

The aggregate production set is $Y = \sum_j Y_j$. It has \emph{no free lunch} if $Y \cap \mathbb{R}^L_+ \subseteq \{0\}$
and is \emph{irreversible} if $Y \cap (-Y) \subseteq \{0\}$.

\begin{lemma}[Inaction, no free lunch and irreversibility]
\label{lem:eq_zero_of_sum_nonneg}
\lean{Econ.ProductionEconomy.aggProd, Econ.ProductionEconomy.eq_zero_of_sum_nonneg}
\leanok
If $0 \in Y_j$ for all $j$ and $Y$ has no free lunch and is irreversible, then $y_j \in Y_j$ and $\sum_j y_j \ge 0$
imply every $y_j = 0$.
\end{lemma}
\begin{proof}
\leanok
No free lunch gives $\sum_j y_j = 0$. For each $k$, $y_k \in Y$ (others at $0$) and $-y_k = \sum_{j \ne k} y_j \in Y$, so
$y_k = 0$ by irreversibility.
\end{proof}

\begin{lemma}[Attainable production is bounded]
\label{lem:bound_attainable}
\lean{Econ.ProductionEconomy.exists_bound_attainable}
\leanok
\uses{lem:eq_zero_of_sum_nonneg}
If each $Y_j$ is closed and convex with $0 \in Y_j$, and $Y$ has no free lunch and is irreversible, then for every $c$ the
set of $(y_j)_j \in \prod_j Y_j$ with $\sum_j y_j \ge c$ is bounded (Debreu 1959, ch.~5).
\end{lemma}
\begin{proof}
\leanok
Otherwise take such $y^n$ with $\|y^n\| > n+1$ and a limit point $d$ of $y^n/\|y^n\|$, so $\|d\| = 1$. Each
$y^n_j/\|y^n\|$ lies on the segment from $0$ to $y^n_j$, hence in $Y_j$, and so $d_j \in Y_j$. Also
$\sum_j y^n_j/\|y^n\| \ge c/\|y^n\| \to 0$, so $\sum_j d_j \ge 0$, and the previous lemma gives $d = 0$, a
contradiction.
\end{proof}

\begin{lemma}[Unique profit maximizers after truncation]
\label{lem:profit_trunc_unique}
\lean{Econ.ProductionEconomy.truncSet, Econ.ProductionEconomy.isClosed_truncSet,
  Econ.ProductionEconomy.isBounded_truncSet, Econ.ProductionEconomy.profitMax_trunc_subsingleton}
\leanok
If $Y_j$ is strictly convex and $p \gg 0$, then $p \cdot y$ has at most one maximizer on $Y_j \cap \{z : \sum_i z_i^2 \le
\rho\}$, a compact set.
\end{lemma}
\begin{proof}
\leanok
The midpoint $m$ of two distinct maximizers lies in the interior of $Y_j$ and satisfies $\sum_i m_i^2 < \rho$ (parallelogram
identity), so $m + tp$ is feasible for small $t > 0$ and has higher profit.
\end{proof}

\begin{theorem}[Existence with unbounded production]
\label{thm:existence_unbounded}
\lean{Econ.ProductionEconomy.exists_walrasianEquilibrium_unbounded}
\leanok
\uses{lem:bound_attainable, lem:profit_trunc_unique, thm:existence_production}
Suppose consumers have regular utilities and $\bar\omega \gg 0$, each $Y_j$ is closed and strictly convex with $0 \in Y_j$,
and $Y = \sum_j Y_j$ has no free lunch and is irreversible. Then a Walrasian equilibrium exists at strictly positive
prices. Production sets need not be bounded. Strict convexity is much stronger than Arrow--Debreu's convexity: it forces
full-dimensional $Y_j$ and excludes constant returns, linear activities and free disposal. This theorem and
Theorem~\ref{thm:existence_production} are not nested.
\end{theorem}
\begin{proof}
\leanok
Bound attainable production by $R$ (with $c = -\bar\omega$) and truncate each $Y_j$ to the ball $\sum_i z_i^2 \le \rho$ with
$\rho = L R^2 + 1$. The truncated economy has compact production sets with unique supply, so it has an equilibrium. Its
production is attainable, hence strictly inside the ball. If some $z \in Y_j$ were more profitable than $y_j$, the points
$y_j + t(z - y_j)$ for small $t > 0$ would lie in the truncated set and be more profitable, a contradiction.
\end{proof}

\begin{theorem}[An unbounded technology]
\label{thm:parabola_economy}
\lean{Econ.GEExamples.parabolaSet, Econ.GEExamples.zero_mem_parabolaSet, Econ.GEExamples.isClosed_parabolaSet,
  Econ.GEExamples.strictConvex_parabolaSet, Econ.GEExamples.noFreeLunch_parabolaSet,
  Econ.GEExamples.irreversible_parabolaSet, Econ.GEExamples.not_isBounded_parabolaSet,
  Econ.GEExamples.parabolaEconomy, Econ.GEExamples.parabolaEconomy_aggProd,
  Econ.GEExamples.parabolaEconomy_exists_equilibrium}
\leanok
\uses{thm:existence_unbounded, thm:sqrt_economy}
The set $\{y : y_1 \le -y_0 - y_0^2\}$ is closed, strictly convex, unbounded, contains $0$, and has no free lunch and is
irreversible. Two square-root consumers owning $(1,0)$ and $(0,1)$, one of whom owns this firm, form an economy with a
Walrasian equilibrium.
\end{theorem}

\section{The Arrow--Debreu theorem with convex production}
\label{sec:existence_arrow_debreu}

\begin{lemma}[Continuous truncated demand]
\label{lem:trunc_demand}
\lean{Econ.truncBudget, Econ.continuous_cost_right, Econ.isClosed_truncBudget, Econ.convex_truncBudget,
  Econ.isClosed_graph_truncBudget, Econ.continuousCorrespondence_truncBudget, Econ.exists_continuous_truncDemand}
\leanok
For $c \ge 0$ and $m > 0$, the truncated budget $\{x \in [0,c]^L : p \cdot x \le \max(w, m)\}$ is a continuous
correspondence of $(p, w)$ with nonempty compact convex values. For continuous strictly quasiconcave $u$, its maximizer
is unique and continuous in $(p, w)$.
\end{lemma}
\begin{proof}
\leanok
Upper hemicontinuity: closed graph inside a fixed box. Lower hemicontinuity: shrink an affordable bundle $y$ to $ty$ with
$t < 1$, which costs strictly less than $\max(w, m) > 0$, a strict inequality that survives perturbation. Uniqueness:
the midpoint of two maximizers is affordable and strictly better. Continuity: Berge.
\end{proof}

\begin{theorem}[Arrow--Debreu existence with convex production]
\label{thm:existence_arrow_debreu}
\lean{Econ.ProductionEconomy.exists_walrasianEquilibrium_arrowDebreu}
\leanok
\uses{thm:existence_arrow_debreu_convex}
Suppose consumers have continuous, strictly quasiconcave, strongly monotone utilities and endowments $\omega_a \gg 0$, each
$Y_j$ is closed and convex with $0 \in Y_j$, and $Y = \sum_j Y_j$ has no free lunch and is irreversible. Then a Walrasian
equilibrium exists at strictly positive prices. Production sets may be unbounded, with constant returns and free
disposal (Arrow--Debreu 1954, Theorem I, here with strictly quasiconcave preferences).
\end{theorem}
\begin{proof}
\leanok
In Lean this is the special case of Theorem~\ref{thm:existence_arrow_debreu_convex}, since strictly quasiconcave
utilities are quasiconcave. The direct argument: bound attainable production by $R$, truncate $Y_j$ to $T_j = Y_j \cap \bar B(0, R+1)$ and consumption to $[0, c]^L$ with
$c > \sum_i \bar\omega_i + |J| R$. Apply Kakutani on $\Delta \times [0,c]^{A \times L} \times \prod_j T_j$ to the
correspondence sending $(p, x, y)$ to the price vectors maximizing $q \cdot z$ over $\Delta$, where
$z = \sum_a x_a - \bar\omega - \sum_j y_j$; the truncated demands at wealth $p \cdot \omega_a + \sum_j \theta_{aj}
\max(0, p \cdot y_j)$ with floor $m_a = \min_i \omega_{ai}$; and the profit maximizers on $T_j$. At a fixed point
profits are nonnegative, wealth is at least $p \cdot \omega_a \ge m_a$, and summing budgets gives $p \cdot z \le 0$, so
$z \le 0$ by price maximization. Hence production is attainable, strictly inside the truncation, and consumption is
strictly inside the box; by convexity firms maximize on $Y_j$, and consumers maximize on their full budget sets (see
the lift-back in Theorem~\ref{thm:existence_arrow_debreu_convex}). Strong monotonicity gives $p \gg 0$ and binding budgets, so $p \cdot z = 0$ and $z = 0$.
\end{proof}

\begin{lemma}[Set-valued truncated demand]
\label{lem:trunc_demand_convex}
\lean{Econ.StrictlyQuasiconcave.quasiconcaveOn, Econ.truncDemand_props}
\leanok
\uses{lem:trunc_demand}
For continuous $u$ quasiconcave on $\mathbb{R}^L_+$, the maximizers of $u$ on the truncated budget set are nonempty and
convex and have a closed graph in $(p, w, x)$. Strictly quasiconcave utilities are quasiconcave on $\mathbb{R}^L_+$.
\end{lemma}
\begin{proof}
\leanok
Berge gives upper hemicontinuity with compact values, hence a closed graph. Two maximizers have equal utility, so by
quasiconcavity every convex combination is at least as good.
\end{proof}

\begin{theorem}[Arrow--Debreu existence with convex preferences]
\label{thm:existence_arrow_debreu_convex}
\lean{Econ.ProductionEconomy.exists_walrasianEquilibrium_arrowDebreu_convex}
\leanok
\uses{lem:bound_attainable, lem:trunc_demand_convex}
Suppose consumers have continuous utilities that are quasiconcave on $\mathbb{R}^L_+$ (convex preferences) and strongly
monotone, and endowments $\omega_a \gg 0$; each $Y_j$ is closed and convex with $0 \in Y_j$; and $Y = \sum_j Y_j$ has no free
lunch and is irreversible. Then a Walrasian equilibrium exists at strictly positive prices.
\end{theorem}
\begin{proof}
\leanok
As before, with the consumer part of the Kakutani correspondence now the set-valued truncated demand. In the lift-back
step, suppose an affordable $x'$ is better than $x_a$. Shrink it to $x'' = \sigma x'$, still better and leaving money
unspent. With the fixed $\delta = (w - p \cdot x'')/(\sum_i p_i + 1)$ the bundle $t x'' + (1-t) x_a + t\delta \mathbf{1}$ is
affordable for every $t \in [0, 1]$ and in the box for small $t > 0$; it is at least as good as $t x'' + (1-t) x_a$, which is at least as good as $x_a$ by quasiconcavity, and is
strictly better by strong monotonicity, contradicting truncated optimality.
\end{proof}

\begin{theorem}[Perfect substitutes]
\label{thm:subst_economy}
\lean{Econ.GEExamples.quasiconcaveOn_uLin, Econ.GEExamples.not_strictlyQuasiconcave_uLin,
  Econ.GEExamples.not_strictlyQuasiconcave_of_prefOf_eq_uLin,
  Econ.GEExamples.stronglyMonotone_uLin, Econ.GEExamples.substEconomy, Econ.GEExamples.substEconomy_aggProd,
  Econ.GEExamples.substEconomy_exists_equilibrium}
\leanok
\uses{thm:existence_arrow_debreu_convex, thm:linear_economy}
The utility $x_0 + x_1$ is continuous, quasiconcave and strongly monotone, and no utility representing it is strictly quasiconcave. Two such consumers
each owning $(1,1)$, with the constant-returns firm of Theorem~\ref{thm:linear_economy}, form an economy with a Walrasian
equilibrium.
\end{theorem}

\begin{theorem}[Constant returns and free disposal]
\label{thm:linear_economy}
\lean{Econ.GEExamples.linearSet, Econ.GEExamples.zero_mem_linearSet, Econ.GEExamples.isClosed_linearSet,
  Econ.GEExamples.convex_linearSet, Econ.GEExamples.smul_mem_linearSet, Econ.GEExamples.freeDisposal_linearSet,
  Econ.GEExamples.noFreeLunch_linearSet, Econ.GEExamples.irreversible_linearSet,
  Econ.GEExamples.linearSet_not_strictConvex, Econ.GEExamples.linearEconomy, Econ.GEExamples.linearEconomy_aggProd,
  Econ.GEExamples.linearEconomy_exists_equilibrium}
\leanok
\uses{thm:existence_arrow_debreu, thm:sqrt_economy}
The cone $\{y : y_0 \le 0,\ y_0 + y_1 \le 0\}$ (convert good 0 into good 1 one for one, dispose freely) is closed,
convex, has constant returns and free disposal, no free lunch, is irreversible and is not strictly convex. Two
square-root consumers each owning $(1,1)$, one of whom owns this firm, form an economy with a Walrasian equilibrium.
\end{theorem}
